% DELTA23 independent mixed Krammer analysis, /root/characters.
% Matrix certificates use F9=F3[i], i^2=-1, integer encoding a+3b.

\subsection{The two mixed Krammer pairs and their exact gentle correlation}

Use the six-dimensional Krammer representation of $B_4$ at
$q=i,t=-i$ over $k=\mathbb F_9$.  Its three literal matrices
$K_1,K_2,K_3$ are constructed in the companion pentagon workstream
from Bigelow's displayed Krammer formula.  Set
\[
 X_1=K_1^2,\quad X_2=K_2^2,\quad X_3=K_3^2,
 \quad X_4=(K_1K_2K_1)^2,\quad X_5=(K_3K_2K_3)^2.
\]
The exact matrices, all the bases below, and the full six-dimensional
intertwiners are recorded in
\path{agents/characters/records/krammer_mixed_correlation.json}; the executable
\path{agents/characters/code/check_krammer_mixed_correlation.py} verifies every
displayed matrix identity in the literal original basis.

In the matrix arrays in this subsection, an integer from zero to
eight means $a+bi\in\mathbb F_9$ when its integer value is $a+3b$,
with $0\le a,b\le2$.  Thus $3=i$, $6=-i$, and $2=-1$.

\paragraph{D23.C9 (proved by exact finite linear algebra and group certificates).}
Each mixed pair $(X_3,X_4)$ and $(X_4,X_5)$ preserves a direct
sum of two three-dimensional spaces.  On the first its image is
a solvable group $H$ of order $108$, with derived orders
\[
 108,27,3,1.
\]
On the second its image is a full unitary group $U(3,3)$ of order
$24192$, with perfect derived group $SU(3,3)$ of order $6048$.
The orders and successive derived orders of the full mixed images are
\[
\begin{array}{c|rrrr}
 \text{pair}&G&G'&G''&G'''=G''''\\\hline
 (X_3,X_4)&1306368&163296&18144&6048\\
 (X_4,X_5)&653184&163296&18144&6048.
\end{array}
\]
Their full linear images inject into the projective images obtained
by quotienting by the scalar $-I$.

There are explicit $C,Q,D\in\mathrm{GL}_6(k)$ such that, writing
$z=(X_4X_5)^{-1}$,
\[
 X_3C=CDz,\qquad X_4C=CX_4,\qquad
 [D,X_4]=[D,X_5]=1,
\]
\[
 zQ=Qz,\qquad X_4Q=QX_5.
\]
Consequently, for every actual multiplier-one gentle word $f$, if
\[
 R=f(X_3,X_4),\qquad r=C^{-1}RC,
\]
then its two mixed values obey the literal six-dimensional identity
\[
 \boxed{f(X_4,X_5)=Q^{-1}rQr^{-1}.}
\]
In particular the second mixed value is not an independent variable
in a finite necessary pentagon test.  The first mixed value belongs
to a group of order $163296$.

\emph{Proof.}
Let $B_{34}$ and $B_{45}$ be the splitting matrices from the
certificate, with the solvable summand first.  Their columns are
the indicated bases of the two invariant spaces, and have rank six.
Writing the corresponding three-dimensional generator pairs as
$(x_H,y_H),(x_U,y_U)$ and $(u_H,v_H),(u_U,v_U)$, the certificate
checks
\[
 B_{34}^{-1}(X_3,X_4)B_{34}
    =((x_H\oplus x_U),(y_H\oplus y_U)),
\]
\[
 B_{45}^{-1}(X_4,X_5)B_{45}
    =((u_H\oplus u_U),(v_H\oplus v_U)).
\]
The invariant-space discovery enumerated all $101$ projective
eigenvectors of one generator in each pair.  Discovery is not
needed as an assumption: the displayed bases and intertwining
identities independently certify the direct sums.

The two solvable block groups were computed as permutation groups
on unions of vector orbits containing a full $\mathbb F_3$-basis
of $k^3$.  Such actions are faithful, since a matrix fixing those
vectors is the identity.  Exact stabilizer chains give orders
$108,27,3,1$ in the derived series.  The first derived group has
exponent three.  The determinant maps the whole group onto
$\mu_4\subset k^\times$, so its kernel, of order $27$, is exactly
the derived group.

For the unitary blocks, the two respective invariant Hermitian
forms are
\[
 J_{34}=\begin{pmatrix}0&1&5\\1&0&3\\8&6&0\end{pmatrix},
 \qquad
 J_{45}=\begin{pmatrix}0&3&6\\6&0&4\\3&7&0\end{pmatrix}.
\]
They are Hermitian for $i\mapsto-i$, have nonzero determinants
$2$ and $1$, and the exact identities $g^*Jg=J$ hold for both
generators.  A faithful union of vector orbits has $224$ vectors;
the stabilizer-chain calculation gives order $24192$ and derived
order $6048$, with the derived group perfect.  The full unitary
group has order
\[
 3^3(3+1)(3^2-1)(3^3+1)=24192.
\]
For a direct count at this small field, there are $252$ vectors of
norm one in dimension three, $24$ in their orthogonal complement
of dimension two, and $4$ in dimension one, giving the same product.
Thus containment and the computed order identify the full group.
The determinant image has order four and the derived group lies
in its kernel; the kernel has order $6048$, so it is exactly the
special unitary group.  This identification does not need a
classification or simplicity theorem for finite groups.

For completeness, the full group orders do not require enumerating
the large six-dimensional image.  Each full mixed image is a
subdirect product of its solvable block $H$ and unitary block $U$.
Any common quotient of $H$ and $U$ is solvable, hence kills the
perfect group $U'=SU(3,3)$ and is abelian.  It therefore also kills
$H'$.  Goursat's lemma gives $H'\times U'$ inside the mixed image;
the image is the full preimage of its determinant-pair image.
For pair $34$, the generator determinant pairs are
\[
 (-1,1),\quad(-i,-i),
\]
generating the equal-parity subgroup of $\mu_4^2$ of order eight.
For pair $45$ they are both $(-i,-i)$ and generate the diagonal
subgroup of order four.  Multiplication by
$|H'|\,|U'|=27\cdot6048$ gives the two stated group orders.
Applying the same solvable/perfect Goursat argument to the derived
subgroup gives $G'=H'\times U'$, and then the remaining orders
follow by taking derived subgroups of a direct product.  Each
solvable block has a common fixed line.  Thus the full mixed image
does not contain scalar $-I$, proving the assertion about the
projective quotient.

Here are small matrices constructing the required full intertwiners:
\[
 C_H=\begin{pmatrix}4&1&0\\3&0&0\\0&0&1\end{pmatrix},
 \quad C_U=\begin{pmatrix}1&3&2\\3&5&0\\0&1&0\end{pmatrix},
\]
\[
 Q_H=\begin{pmatrix}4&0&6\\2&1&0\\1&0&0\end{pmatrix},
 \quad Q_U=\begin{pmatrix}2&0&2\\2&5&0\\6&2&1\end{pmatrix}.
\]
Put $z_H=(u_Hv_H)^{-1}$ and $z_U=(u_Uv_U)^{-1}$.  Direct
multiplication gives
\[
 x_HC_H=C_Hz_H^3,\quad y_HC_H=C_Hu_H,
 \quad x_UC_U=C_U(-z_U),\quad y_UC_U=C_Uu_U,
\]
and $z_H^2$ commutes with $u_H,v_H$.  For each block one also has
$z_jQ_j=Q_jz_j$ and $u_jQ_j=Q_jv_j$.  Define
\[
 \begin{split}
 C&=B_{34}(C_H\oplus C_U)B_{45}^{-1},\\
 Q&=B_{45}(Q_H\oplus Q_U)B_{45}^{-1},\\
 D&=B_{45}(z_H^2\oplus(-I_3))B_{45}^{-1}.
 \end{split}
\]
All three matrices are invertible; the stated literal
six-dimensional identities follow by conjugating the block
identities.  The companion computation also checks them directly
in the original basis, independently of this deduction.

If a profinite word $f$ belongs to $F'$, multiplying an input by
an element centralizing both inputs does not change its value:
the extra factor is raised to the corresponding abelian exponent,
which is zero.  Consequently
\[
 C^{-1}f(X_3,X_4)C=f(Dz,X_4)=f(z,X_4)=r.
\]
The other intertwiner gives
$f(z,X_5)=Q^{-1}rQ$.  The multiplier-one two-cycle equation says
$f(X_5,z)=f(z,X_5)^{-1}=Q^{-1}r^{-1}Q$.  In the source's exact
order, its three-cycle equation is
\[
 f(z,X_4)f(X_5,z)f(X_4,X_5)=I
\]
(see DELTA\_ON\_GT P.9, the displayed two- and three-cycle
equations in the normalized lift proof).  Substitution and inversion
give $f(X_4,X_5)=Q^{-1}rQr^{-1}$, as asserted.  Finally
$f\in F'$ puts its first mixed value in $G_{34}'$, whose order
has already been computed.\hfill$\square$

\paragraph{D23.C10 (proved; the nilpotent residual does not erase this block).}
Under either mixed-pair evaluation, the image of $F''$ is
$H''\times SU(3,3)$, of order $3\cdot6048$.  The image of
$R_{\rm nil}(F'')$ is exactly
\[
 \{1\}\times SU(3,3).
\]
These assertions concern the images of all words in the indicated
subgroups.  They do not assert that actual gentle words fill those
images or satisfy the pentagon there.

\emph{Proof.}
Surjectivity of the free evaluation and D23.C9 give the claimed
image of $F''$ as the second derived subgroup.  The $H''$ factor
has order three, so it is a finite nilpotent quotient and is
annihilated by $R_{\rm nil}(F'')$.  Let $N$ be the image of that
residual in the perfect $SU(3,3)$ factor.  The finite group
$SU(3,3)/N$ is a quotient of the pronilpotent group
$F''/R_{\rm nil}(F'')$, hence is nilpotent.  It is also perfect,
being a quotient of a perfect group.  A finite perfect nilpotent
group is trivial, so $N=SU(3,3)$.\hfill$\square$
